Two-scale Lorenz-96 is an idealised toy system, and the ranking of closures found here may not transfer to real atmosphere or ocean models. This is a small CPU study of one 8-variable configuration ($K{=}8$, $J{=}32$, $F{=}20$) with two values of $c$, five seeds in the main comparison and three in ablations; the $p$-values are uncorrected Welch tests on $n\le5$ and the equivalence claim in H3 rests on non-significance. All baselines are our reimplementations with fixed, untuned hyperparameters; a tuned MLP or a different closure family could change the ranking. All closures depend only on the local slow variable in the main comparison and are trained offline: we do not test coupled/online training, which the cited work proposes as a remedy for instabilities and biases~\cite{rasp2019coupled}, and ``stable'' means stable over 20 runs of 400 time units. The reference climate is itself a finite run, and the independent-truth floor shows how large that sampling error is (Table~\ref{tab:floor}). The AR(1) noise is independent between sectors and uninitialised at forecast start, so its forecast skill reflects an ensemble mean with 10 members. The ablations reuse the test protocol and the forcing shift is mild and one-dimensional; the degree-2 failure and the noise-amplitude effects are not explained mechanistically. Two runs exceeded the intended 6-minute wall time (Section~\ref{sec:setup}). The step-size pre-check is a single run that compares only the mean and variance of $X$. A full study would add more seeds, tuned neural models, state-dependent and multiplicative noise, online training, ensemble-skill scores (CRPS) and larger systems.
